\documentclass[12pt,a4paper]{article}

% Drake_preamble.tex - LaTeX Preamble Template
% 作者: 謝慕揚, MD, PhD, FESC

% Including essential packages for Chinese typesetting
\usepackage{xeCJK}
\usepackage{fontspec}
% Configuring Chinese font with Noto Serif CJK TC, fallback to SimSun
\IfFontExistsTF{Noto Serif CJK TC}{
  \setCJKmainfont{Noto Serif CJK TC}
  \setCJKsansfont{Noto Serif CJK TC}
  \setCJKmonofont{Noto Serif CJK TC}
}{\setCJKmainfont{SimSun}
  \setCJKsansfont{SimSun}
  \setCJKmonofont{SimSun}
}

% Including other necessary packages
\usepackage{geometry}
\usepackage{titlesec}
\usepackage{enumitem}
\usepackage{xcolor}
\usepackage{colortbl}
\usepackage{tcolorbox}
\usepackage{booktabs}
\usepackage{longtable}
\usepackage{array}
\usepackage{graphicx}
\usepackage{tikz}
\usepackage{fancyhdr}
\usepackage{multicol}
\usepackage{datetime2}
\usepackage{hyperref}
\usepackage{mdframed}
\usepackage{amsmath}
\usepackage{amssymb}
\usepackage{multirow}

% Defining colors
\definecolor{emergency_red}{RGB}{186,24,27}
\definecolor{emergency_blue}{RGB}{0,114,188}
\definecolor{emergency_gray}{RGB}{120,120,120}
\definecolor{highlight_yellow}{RGB}{255,250,205}

\hypersetup{
    colorlinks=true,
    linkcolor=emergency_blue,
    citecolor=emergency_blue,
    urlcolor=emergency_blue,
    pdfborder={0 0 0}
}

% Defining page geometry
\geometry{
  top=2.5cm,
  bottom=2.5cm,
  left=2cm,
  right=2cm,
  headheight=25pt
}

% Adjusting topmargin to compensate for increased headheight
\addtolength{\topmargin}{-10pt}

% Setting title formats
\titleformat{\section}[block]{\Large\bfseries\color{emergency_red}}{\thesection}{10pt}{}
\titleformat{\subsection}[block]{\large\bfseries\color{emergency_blue}}{\thesubsection}{8pt}{}
\titleformat{\subsubsection}[block]{\normalsize\bfseries\color{emergency_gray}}{\thesubsubsection}{6pt}{}

% Defining custom environment for important notes
\newmdenv[
    linecolor=emergency_red,
    backgroundcolor=highlight_yellow,
    linewidth=2pt,
    topline=true,
    bottomline=true,
    leftline=true,
    rightline=true,
    innertopmargin=10pt,
    innerbottommargin=10pt,
    innerleftmargin=10pt,
    innerrightmargin=10pt
]{importantbox}

% Defining note command with tcolorbox
\newcommand{\note}[1]{
    \par
    \noindent
    \begin{tcolorbox}[
        colback=emergency_blue!5!white,
        colframe=emergency_blue,
        title=\textbf{重點提醒},
        fonttitle=\bfseries,
        boxrule=0.5pt,
        arc=3pt,
        left=6pt,
        right=6pt,
        top=6pt,
        bottom=6pt
    ]
    #1
    \end{tcolorbox}
}

% Key findings box
\newcommand{\keyfinding}[1]{
    \par
    \noindent
    \begin{tcolorbox}[
        colback=yellow!10!white,
        colframe=orange!75!black,
        title=\textbf{核心發現摘要},
        fonttitle=\bfseries,
        boxrule=0.5pt,
        arc=3pt
    ]
    #1
    \end{tcolorbox}
}

% Defining custom date format
\newcommand{\mydate}{\the\year-\the\month-\the\day}

\begin{document}

%% Title Page
\begin{center}
{\LARGE\bfseries\textcolor{emergency_red}{JACC Behind the Science Podcast}}\\[0.5cm]
{\Large\textbf{Loss of Y Chromosome and Major Cardiovascular Events in Older Men}}\\[0.3cm]
{\large Y 染色體缺失與老年男性重大心血管事件：ASPREE 試驗分析}\\[0.5cm]
{\normalsize JACC. 2026 January Behind the Science Podcast}\\[0.3cm]
{\normalsize 整理：謝慕揚 MD, PhD, FESC \quad 日期：\mydate}
\end{center}

\vspace{0.5cm}

\keyfinding{
\begin{itemize}[leftmargin=*]
    \item Loss of Y chromosome (LOY) 是指循環 leukocytes 完全失去 Y 染色體，這在 70 歲以上男性中相當常見（20-40\%）
    \item 在 ASPREE 試驗中，baseline 時的 LOY 可前瞻性預測未來 myocardial infarction (MI) 風險
    \item 關鍵發現：LOY 與 MI 顯著相關，但與 ischemic stroke 無關——可能反映不同的病理機制
    \item Dose-response 效應：$>$23\% cells 有 LOY 者，non-fatal MI 的 HR = 1.68（95\% CI 1.16-2.45）
    \item LOY 與 CHIP (clonal hematopoiesis of indeterminate potential) 呈反向相關，可能是 competing processes
\end{itemize}
}

\section{研究背景}

\subsection{什麼是 Loss of Y Chromosome (LOY)?}

\begin{itemize}
    \item \textbf{定義}：循環 leukocytes（白血球）完全失去整條 Y 染色體的現象
    \item \textbf{發生率}：70 歲以上男性中，20-40\% 會出現此現象
    \item \textbf{程度}：受影響的 leukocytes 比例可從 10\% 到 50\% 不等
    \item \textbf{發現者}：約 10-15 年前由瑞典 Uppsala 大學的 Lars Forsberg 教授發現
\end{itemize}

\subsection{LOY 與疾病的關聯}

過去研究顯示 LOY 與多種疾病相關：
\begin{itemize}
    \item \textbf{癌症}：血液癌症及實體腫瘤
    \item \textbf{神經退化}：Alzheimer's disease
    \item \textbf{眼科疾病}：Macular degeneration（黃斑部病變）
    \item \textbf{心血管疾病}：但過去研究多為已有 CV disease 或 CKD 的患者
\end{itemize}

\note{
本研究的新貢獻：首次在「健康」的老年男性中，前瞻性地證明 LOY 可預測未來心血管事件風險。
}

\section{ASPREE 試驗簡介}

\subsection{試驗設計}

\begin{longtable}{p{4cm}p{10cm}}
\toprule
\textbf{項目} & \textbf{內容} \\
\midrule
\endfirsthead
\toprule
\textbf{項目} & \textbf{內容} \\
\midrule
\endhead
試驗名稱 & ASPREE (ASPirin in Reducing Events in the Elderly) \\
\midrule
設計 & 大型隨機對照試驗 \\
\midrule
介入 & Low-dose aspirin vs placebo \\
\midrule
Primary endpoint & Disability-free survival \\
\midrule
參與者 & 19,000 人（澳洲及美國） \\
\midrule
年齡 & $\geq$65 歲（少數族裔）或 $\geq$70 歲（多數）\newline 平均年齡約 74 歲 \\
\midrule
關鍵特點 & 所有參與者在 enrollment 時\textbf{無心血管疾病病史}\newline 由 GP 進行臨床評估確認 \\
\midrule
Biobank & 約 75\% 參與者提供 baseline 血液樣本 \\
\bottomrule
\end{longtable}

\subsection{為何 ASPREE 適合研究 LOY?}

\begin{enumerate}
    \item \textbf{Primary prevention setting}：參與者在 baseline 時無 CV disease
    \item \textbf{老年族群}：平均 74 歲，正是 LOY 開始顯著增加的年齡
    \item \textbf{大規模 biobank}：可進行基因分析及 biomarker 測量
    \item \textbf{臨床 adjudicated endpoints}：心血管事件經專家確認
\end{enumerate}

\section{主要研究發現}

\subsection{LOY 與心血管事件}

\begin{itemize}
    \item LOY 與 MACE（major adverse cardiovascular events）及其組成部分相關
    \item \textbf{關鍵發現}：LOY 主要與 myocardial infarction 相關，\textbf{但與 ischemic stroke 無關}
\end{itemize}

\subsection{Dose-Response 效應}

研究將 LOY 程度分為三類：
\begin{longtable}{p{4cm}p{5cm}p{5cm}}
\toprule
\textbf{LOY 程度} & \textbf{定義} & \textbf{說明} \\
\midrule
\endfirsthead
\toprule
\textbf{LOY 程度} & \textbf{定義} & \textbf{說明} \\
\midrule
\endhead
低度或無 LOY & $<$5\% cells 失去 Y 染色體 & 可能為 noise/artifact \\
\midrule
中度 LOY & 5-22\% cells 失去 Y 染色體 & 介於 signal 與 noise 之間 \\
\midrule
高度 LOY (Substantial) & $>$23\% cells 失去 Y 染色體 & 明確的 pathological signal \\
\bottomrule
\end{longtable}

\subsection{主要結果}

\begin{itemize}
    \item \textbf{Non-fatal MI}（高度 LOY vs 低度 LOY）：
    \begin{itemize}
        \item Fully adjusted HR = 1.68（95\% CI 1.16-2.45）
        \item 調整 CHIP 後 HR = 1.56（仍顯著）
        \item 排除 CHIP 患者後 HR = 1.9（信號增強）
    \end{itemize}
    \item \textbf{Ischemic stroke}：95\% CI 跨越 1（無顯著相關）
    \item \textbf{MACE}：95\% CI 跨越 1（無顯著相關）
\end{itemize}

\note{
為何 LOY 與 MI 相關但與 stroke 無關？可能的解釋：
\begin{itemize}
    \item MI 主要來自 plaque rupture 後的局部 inflammation
    \item Stroke 多為來自大動脈或心臟的 emboli（如 atrial fibrillation）
    \item 類似的 exposure-outcome 差異也見於：Lipoprotein(a)、LDL、CETP gene——均與 MI 相關但與 stroke 無關
\end{itemize}
}

\section{外部驗證：UK Biobank}

\subsection{驗證結果}

\begin{itemize}
    \item 在 UK Biobank 中複製了 LOY 與 MI/MACE 的關聯
    \item 效應量較小（可能原因見下）
    \item Stroke 在 UK Biobank 中同樣無顯著關聯
\end{itemize}

\subsection{ASPREE vs UK Biobank 差異}

\begin{longtable}{p{4cm}p{5cm}p{5cm}}
\toprule
\textbf{項目} & \textbf{ASPREE} & \textbf{UK Biobank} \\
\midrule
\endfirsthead
\toprule
\textbf{項目} & \textbf{ASPREE} & \textbf{UK Biobank} \\
\midrule
\endhead
年齡 & 平均 74 歲 & 平均 57 歲 \\
\midrule
LOY 盛行率 & 較高（老年族群） & 較低（LOY 在 65+ 才開始顯著） \\
\midrule
CV disease 排除 & 臨床評估確認 & ICD-10 codes（較不精確） \\
\midrule
Endpoint 確認 & 臨床 adjudication & ICD-10 codes \\
\midrule
Selection bias & 較小 & 較大（9M 邀請信 → 500K 參與） \\
\bottomrule
\end{longtable}

\section{LOY 與 CHIP 的關係}

\subsection{CHIP 簡介}

\textbf{CHIP (Clonal Hematopoiesis of Indeterminate Potential)}：
\begin{itemize}
    \item 造血幹細胞獲得體細胞突變並 clonal expansion
    \item 也是 age-related 現象
    \item 與心血管疾病風險增加相關
\end{itemize}

\subsection{LOY 與 CHIP 的關係}

\begin{itemize}
    \item 原始假設：LOY + CHIP = 「Double whammy」，風險更高
    \item \textbf{實際發現}：LOY 與 CHIP 呈\textbf{反向相關}（inversely correlated）
    \begin{itemize}
        \item 有 LOY 者較不易有 CHIP
        \item 有 CHIP 者較不易有 LOY
    \end{itemize}
    \item 推論：LOY 與 CHIP 可能是\textbf{獨立或 competing processes}
    \item 可能發生在不同的 blood cell subsets
\end{itemize}

\section{LOY 測量方法}

\subsection{Microarray 方法（Gold Standard）}

\begin{enumerate}
    \item 計算所有 non-sex chromosomes 的 SNP microarray intensity 平均值
    \item 計算 Y chromosome 所有 probes 的 intensity 平均值
    \item 比較兩者的比值（ratio）估計 LOY
    \item 轉換為「失去 Y 染色體的細胞百分比」
\end{enumerate}

\subsection{訊號與雜訊}

\begin{itemize}
    \item 低端測量有 signal-to-noise 問題
    \item 部分數據可能出現「正值」（看似有更多 Y 染色體）——這是 artifact
    \item $>$23\% LOY 可確定為真實的 pathological signal
\end{itemize}

\subsection{Whole Genome Sequencing 驗證}

\begin{itemize}
    \item 在約 1,200-1,500 名男性亞群中進行
    \item 透過 Y chromosome 的 sequencing depth 估計 LOY
    \item 與 microarray 結果高度相關（r = 0.98）
    \item Microarray 在低端可能比 sequencing 更有雜訊
\end{itemize}

\section{可能的機制}

\subsection{動物實驗證據}

Sano 等人（約 2 年前）的研究顯示：
\begin{itemize}
    \item LOY 在動物模型中導致 fibrosis
    \item \textbf{心肌}是最受影響的肌肉組織
    \item 可能機制：MI 後 remodeling
\end{itemize}

\subsection{推測的機制}

\begin{itemize}
    \item \textbf{Pro-fibrotic signaling}：TGF-$\beta$（transforming growth factor beta）
    \item \textbf{Macrophage 功能}：LOY macrophages 可能促進 fibrosis
    \item \textbf{Y 染色體上的基因}：可能影響 blood pressure 及 cholesterol（已知 MI 風險因子）
\end{itemize}

\section{臨床意義與未來方向}

\subsection{目前限制}

\begin{itemize}
    \item LOY 檢測目前\textbf{尚未在臨床實驗室常規提供}
    \item 仍是「emerging risk factor」
    \item 是否為 causative 或 byproduct 仍需更多研究
\end{itemize}

\subsection{未來研究方向}

\begin{enumerate}
    \item \textbf{風險預測}：將 LOY 納入現有心血管風險評分（如 SCORE-OP）
    \item \textbf{治療介入}：是否可預防或逆轉 LOY？
    \item \textbf{機制研究}：LOY 是否發生在特定 cell subsets？
    \item \textbf{其他族群驗證}：不同人種、不同年齡層
\end{enumerate}

\note{
重要觀點：LOY 可能是「age-related genomic changes in blood」的一部分，這類變化在心血管疾病風險評估中越來越重要。本研究也特別感謝 ASPREE 試驗 PI John McNeil 教授建立的 biobank 基礎建設。
}

\section{結論}

\begin{enumerate}
    \item LOY 是老年男性中常見的 age-related 現象（70 歲以上 20-40\%）
    \item 在健康老年男性中，LOY 可前瞻性預測 MI 風險
    \item LOY 與 MI 的關聯獨立於傳統心血管風險因子及 CHIP
    \item LOY 與 stroke 無關，可能反映 MI 與 stroke 不同的病理機制
    \item 未來需要更多研究確定 LOY 的臨床應用價值
\end{enumerate}

\section{參考資料}

\begin{enumerate}
    \item Lacaze P, Manierato A, et al. Loss of Y Chromosome and Major Cardiovascular Events in a Prospective Study of Older Men. \textit{J Am Coll Cardiol}. 2026.
    \item Sano S, et al. Hematopoietic loss of Y chromosome leads to cardiac fibrosis and heart failure mortality. \href{https://doi.org/10.1126/science.abn3100}{\textit{Science}. 2022;377:292-297.}
    \item Forsberg LA. Loss of chromosome Y (LOY) in blood cells is associated with increased risk for disease and mortality in aging men. \href{https://doi.org/10.1007/s00439-017-1799-2}{\textit{Hum Genet}. 2017;136:657-663.}
\end{enumerate}

\end{document}
